\chapter*{Abstract}
\addcontentsline{toc}{chapter}{Abstract}
\markboth{Abstract}{Abstract}

This report presents \rat{}, an integrated static and dynamic analysis platform for detecting
\emph{Remote Access Trojans} (RATs) in Windows executables (and C\# source with reduced scope),
developed as the final-year project in Computer Science Engineering at \ubi{}. The solution
combines a Python backend (FastAPI) with a seven-phase pipeline (PE metadata, ILSpy, heuristics,
deobfuscation, YARA, risk scoring, native/Ghidra when needed, and reporting), a React web
interface (with optional Gemini assistant), a WPF desktop application, and an isolated
Hyper-V sandbox with two independent orchestration paths (HTTP vm-agent and PowerShell Direct
with SHA-256 integrity checks). Validation relies on 220 automated tests, proxy detection
metrics on 59 labelled synthetic profiles stratified across five severity levels
(precision/recall/F1 and level accuracy 1.00), a benign corpus of eight compiled
PEs analysed statically (specificity 1.00, zero false positives), public MalwareBazaar IOC metadata
(without host-side malware storage), a formal static benchmark (median 36.1~s, $n=5$), synthetic scoring profiles aligned with
\filepath{test\_risk\_scoring.py}, and benign execution via \filepath{benign-vm-test.exe} on
Path~B. The system was designed for explainability and local isolation rather than commercial cloud
sandboxes, targeting academic and laboratory use.

\Needspace{4\baselineskip}
\vspace{0.8em}
\noindent\textbf{Keywords:} malware analysis; RAT; sandbox; Hyper-V; YARA; FastAPI;
software engineering; cybersecurity.
